\section{引言}
\label{sec:introduction}

洪涝、地震与野火之后，指挥员更希望对无人机说``检查道路北侧那栋完全损毁建筑''，而不是给出精确坐标。
执行这类指令需要语言条件目标搜索、导航与损伤评估，并且必须承认：第一次俯视观测常常不够。
鸟瞰条件下，小目标、配准误差、类别极不平衡以及卫星--低空域偏移都会使一次性置信判断不可靠。

视觉语言导航（VLN）多在第一人称三维环境中研究\cite{anderson2018vision,liu2023aerialvln}。
灾害测绘则广泛使用地理配准的星载或航拍影像与灾前/灾后配对，如 xBD/xView2\cite{gupta2019xbd}。
空中语言导航近期进一步引入分层语义规划与拓扑记忆\cite{zhang2025citynavagent,gao2024stmr,shah2023lmnav}。
本文连接上述线索，构建一套面向灾后侦察的语言条件智能体，并在该系统上检验校准不确定性驱动的主动复检。
验证出口刻意分层：
把``校准不确定性能否转化为更好的观测分配''放在离线、无导航的预算实验中回答；
把``现有栈在哪一层失败''放在 oracle 阶梯中诊断；
把完整飞行闭环作为现实检验；
把系统本身作为可运行、可复现的工程载体单独陈述（\secref{sec:architecture}），不把模块存在性写成任务增益。
中心机制的形式化（动作集、归一化熵、期望熵减、共形触发）见 \secref{sec:baro}。

本文贡献如下。
\begin{itemize}
\item \textbf{智能体系统与仿真/监控平台}（工程存在性声明）：双模式编排器（一次性任务规划与闭环语言导航）、分层语义规划与记忆栈、地理配准 xBD 俯视仿真，以及实时人机协同控制台。
这些模块已实现并可复现（\secref{sec:architecture}）。该声明不蕴含导航成功率或复检增益；尤其禁止把定位 Dice、提议召回或记忆图的存在写成判断/到达改善。
\item 面向航空应急侦察的任务定义：把校准不确定性转为预算受限的下降--居中动作，并以有效 GSD 作为物理可审计量。
\item 离线观测分配实验：在事件不相交裁块上比较校准/未校准/共形/oracle 参照。当前分类器下预测几乎不随 GSD 翻转，机制主张在该数据上不可识别。
\item \textbf{对该负结果的反事实排除}：训练/验证曲线核查排除``分类器从未训练成功''，类别加权重训排除``完全未处理类别不平衡''，
标准切分基线对标把``评测协议过严''与``架构/训练规模不足''两个可能原因分离——这一组诊断实验本身构成把负结果与实现缺陷区分开的可复用方法论。
\item 泄漏安全协议与诊断：事件级切分、功效分析、oracle 阶梯瓶颈归因（指认层，附命名修正）、以及 RescueNet 低空评估域。
端到端历史全零成功的三处实现缺陷排查记录在 \secref{sec:x0x3}，作为可复现性说明而非独立科学贡献。
\end{itemize}
